\scanpage{24}
\uline{Corollaires} ``purement géométriques'':

\uline{Corollaire} 1 Soient $X,Y$ des schémas \inserted{de t.f. réduits connexes} sur un corps alg. clos, $k$, et $f,g:X\longrightarrow Y$ deux morphismes, on suppose que $\pi_1(f),\pi_1(g):\pi_1(X)\rightrightarrows\pi_1(Y)$ sont égaux (en tant qu'hom. extérieurs). Alors

a) Si $Y$ \deleted{est un \unsure{schéma de groupe}} \inserted{se plonge par un $i:Y\longrightarrow G$ en un groupe algébrique commutatif \unsure{$G$}} extension d'une VA par un tore, il existe un $a\in Y$ (unique) tel que $g(x)=f(x)+a$ $\forall x\in X(k)$, i.e. $(i\circ g)=\tau_a\circ(i\circ f)$ ($\tau_a$ la translation par $a$).

b) Si $Y$ est une variété élémentaire d'\unsure{Artin} ou \deleted{\unsure{une}} \unsure{fibrée successive des} \unsure{variétés} élémentaires, \inserted{et $X$ \deleted{normal} \inserted{irred.} et si $f$ ou $g$ est dominant,} alors $f=g$.

\margin{\uuline{b)} \deleted{Ce cas est \unsure{prouvé} \unsure{par nous} en car. $0$ \unsure{pour les courbes}; en car. $p>0$, on ne sait \unsure{rien sur la} structure du $\pi_1$ (\unsure{pro-$p$}, \unsure{pro-prime à $p$}), \unsure{et même} les centres des $\pi_1$\ldots}}
\margin{\unsure{Une conjecture forte établirait l'unicité des} \unsure{morphismes}, avec $X$ irréductible et \unsure{quels que soient les} \unsure{applications} $f,g$, \unsure{sans les supposer|faut les supposer} \unsure{dominantes}.}
\margin{En fait, dans a) il suffit de supposer que $H_1(f,\mathbb Z_\ell)=H_1(g,\mathbb Z_\ell)$ pour un $\ell$ premier à la car. \deleted{\unsure{exp.}}}

\uline{Démonstration} a) L'unicité de $a$ est claire -- et \deleted{\unsure{il}} \inserted{l'existence} revient, \inserted{remplaçant $Y$ par $G$ et} en choisissant une ``origine'' $a\in X(k)$, et en se ramenant au cas $f(a)=g(a)=0$, à prouver qu'alors $f=g$. D'ailleurs l'hyp. $\pi_1(f)=\pi_1(g)$ peut être affaiblie en l'hyp. que \deleted{$f$ est} $H_1(f,\mathbb Z_\ell)=H_1(g,\mathbb Z_\ell)$, \inserted{$\ell$ premier à car.,} car ici il suffit
\margin{23}
